% Part: proof-theory
% Chapter: sequent-calculus
% Section: rules-G3c

\documentclass[../../../include/open-logic-section]{subfiles}

\begin{document}

\begin{table}
\begin{tabular}{@{}c@{}c@{}}
\multicolumn{2}{@{}c@{}}{સ્વયંસિદ્ધો}\\[2ex]
$!C, \Gamma \Sequent \Delta, !C$ & $\lfalse, \Gamma \Sequent \Delta$
\\[2ex]
\multicolumn{2}{@{}c@{}}{તાર્કિક નિયમો}\\[2ex]
\Axiom$\lnot !A, \Gamma \fCenter \Delta, !A $
\RightLabel{\LeftR{\lnot}}
\UnaryInf$\lnot !A, \Gamma \fCenter \Delta$
\DisplayProof
&
\Axiom$!A, \Gamma \fCenter$
\RightLabel{\RightR{\lnot}}
\UnaryInf$ \Gamma \fCenter \Delta, \lnot !A $
\DisplayProof
\\[4ex]
\Axiom$ !A, !B, \Gamma \fCenter \Delta$
\RightLabel{\LeftR{\land}}
\UnaryInf$ !A \land !B, \Gamma \fCenter \Delta$
\DisplayProof
&
\Axiom$\Gamma \fCenter \Delta, !A$
\Axiom$ \Gamma \fCenter \Delta, !B$
\RightLabel{\RightR{\land}}
\BinaryInf$ \Gamma \fCenter \Delta, !A \land !B $
\DisplayProof
\\[4ex]\Axiom$!A, \Gamma \fCenter \Delta$
\Axiom$!B, \Gamma \fCenter \Delta$
\RightLabel{\LeftR{\lor}}
\BinaryInf$!A \lor !B, \Gamma \fCenter \Delta$
\DisplayProof
&
\Axiom$\Gamma \fCenter \Delta, !A, !B$
\RightLabel{\RightR{\lor}}
\UnaryInf$ \Gamma \fCenter \Delta, !A \lor !B$
\DisplayProof
\\[4ex]
\Axiom$!A \lif !B, \Gamma \fCenter \Delta, !A$
\Axiom$ !B, \Gamma \fCenter \Delta$
\RightLabel{\LeftR{\lif}}
\BinaryInf$!A \lif !B, \Gamma \fCenter \Delta$
\DisplayProof
&
\Axiom$ !A, \Gamma \fCenter !B$
\RightLabel{\RightR{\lif}}
\UnaryInf$ \Gamma \fCenter \Delta, !A \lif !B $
\DisplayProof
\\[4ex]
\Axiom$ !A(t), \lforall[x][!A(x)], \Gamma \fCenter \Delta$
\RightLabel{\LeftR{\lforall}}
\UnaryInf$ \lforall[x][!A(x)],\Gamma \fCenter \Delta$
\DisplayProof
&
\Axiom$ \Gamma \fCenter !A(a) $
\RightLabel{\RightR{\lforall}}
\UnaryInf$ \Gamma \fCenter \lforall[x][!A(x)]$
\DisplayProof
\\[4ex]
\Axiom$ !A(a), \Gamma \fCenter \Delta $
\RightLabel{\LeftR{\lexists}}
\UnaryInf$ \lexists[x][!A(x)], \Gamma \fCenter \Delta$
\DisplayProof
&
\Axiom$ \Gamma \fCenter \Delta, \lexists[x][!A(x)], !A(t) $
\RightLabel{\RightR{\lexists}}
\UnaryInf$ \Gamma \fCenter \Delta, \lexists[x][!A(x)]$
\DisplayProof
\end{tabular}
\caption{અંતઃપ્રજ્ઞાવાદી તર્ક માટે અનેક ફલિત ધરાવતા \Log{mG3i}ના નિયમો.
$\RightR{\forall}$ અને $\LeftR{\exists}$માં !!{constant}~$a$
ફલિત-સિક્વન્ટમાં આવવો નહીં. $\Gamma$ અને $\Delta$ એ
!!{formula}ના બહુગણ છે. સ્વયંસિદ્ધોમાં !!{formula}~$!C$
પરમાણ્વીય હોવું જોઈએ. મૂળ કૅપ્શન \Log{mG1m}ને
$\lfalse,
\Gamma \Sequent \Delta$ સ્વરૂપનાં સ્વયંસિદ્ધો વિનાના \Log{mG1i}
તરીકે પણ વર્ણવે છે.}
\ollabel{tab:G3c}
\end{table}
\sourcecorrection{OLMD-337}{મૂળ સર્વપરિમાણકના ડાબા નિયમના આધારમાં મુખ્ય સૂત્ર અને સંદર્ભ વચ્ચેનો અલ્પવિરામ ઉમેર્યો છે.}
\sourcecorrection{OLMD-338}{મૂળ કૅપ્શનના અંતે mG1i અને mG1m આવે છે, પરંતુ આ અધ્યાય તેમના નિયમો આપતો નથી. તેથી આ વાક્યને મૂળનો ઉલ્લેખ માની જાળવ્યું છે. તેને mG3iમાંથી ન્યૂનતમ તંત્ર બનાવવાની વ્યાખ્યા ગણવી નહીં; એવા તંત્રના નિયમો માટે જરૂરી માહિતી અહીં મળતી નથી.}

\end{document}
